Pure physiological table
of Universal Principles of Natural Science

1.
Axioms
of intuition

2.
Anticipations
of perception

<Axioms of Intuition.>

Principle of pure understanding: All appearances are, as
regards their intuition, extensive magnitudes.

<Their principle is: All intuitions are
extensive magnitudes.>

<Proof
All appearances contain, as regards their form, an intuition in space
and time, which grounds all of them a priori. They cannot be appre-
hended, therefore, i.e., taken up into empirical consciousness, except
through the synthesis of the manifold through which the representa-
tions of a determinate space or time are generated, i.e., through the
composition of that which is homogeneous and the consciousness of
the synthetic unity of this manifold (of the homogeneous). Now the
consciousness of the homogeneous manifold in intuition in general, inso-
far as through it the representation of an object first becomes possible,
is the concept of a magnitude (quanti). Thus even the perception of an
object, as appearance, is possible only through the same synthetic unity
of the manifold of given sensible intuition through which the unity of
the composition of the homogeneous manifold is thought in the con-
cept of a magnitude, i.e., the appearances are all magnitudes, and in-
deed extensive magnitudes, since as intuitions in space or time they
must be represented through the same synthesis as that through which
space and time in general are determined.>

I call an extensive magnitude that in which the representation of the
parts makes possible the representation of the whole (and therefore
necessarily precedes the latter).


